\section{Results}\label{sec:results}
The complete data ($6$ angles $\times$ $4$ meshes $\times$ $7$ degrees $\times$ $2$ elements
$\times$ $3$ problems $=1008$ solutions) are in \texttt{results/raw/study\_a*.csv}.
Rates are given for the finest pair $N=4\to8$; ``--'' marks errors at the round-off
level (for \PM\ and $p\ge5$: less than ten times the error of the polynomial field on
the same mesh, which is the empirical round-off level, \cref{sec:poly}); $^\dagger$
marks angles with folded elements on the meshes involved (\cref{tab:folded}).

\subsection{Reproduction of Section 5.9 for $p=2$}\label{sec:p2}
For $p=2$ and the meshes $N=1,2,4$ of the paper, the least-squares slopes of $\eE$
(Lam\'e) are $1.90$, $1.95$, $1.14$, $1.57$ for \PM\ and $1.68$, $1.11$, $1.07$, $0.69$
for \PP\ at $\alpha=0^\circ,30^\circ,45^\circ,60^\circ$
(\citet[Fig.~20]{Rajendran2010technique}: PM-QUAD8$^\ast$ $2.1$, $1.85$, $1.42$, $1.21$;
QUAD8 $1.9$, $0.95$, $0.94$, $0.98$). The trends agree: the conventional element drops
to first order on distorted meshes, the unsymmetric element stays close to second order
for $\alpha\le30^\circ$ and becomes erratic beyond. Exact agreement is not expected,
since the present $p=2$ element is the 9-node Lagrange element with $Q_2$ metric
functions in skew coordinates, not the 8-node serendipity element of the paper. On
the finer pair $N=4\to8$, \PM\ ($p=2$) recovers rate $\approx2$ for all
$\alpha\le45^\circ$ (\cref{tab:rates-lame-E,tab:rates-lame-N}).

\subsection{Polynomial field: distortion immunity}\label{sec:poly}
For $p\ge5$ the degree-5 field lies in the \PM\ trial space on every element.
\Cref{tab:poly} shows the nodal error: on meshes with $\det\bm J>0$ and moderate
distortion ($\alpha\le30^\circ$) it is $10^{-14}$--$10^{-8}$, i.e.\ the field is
reproduced up to round-off on every mesh, as predicted. The round-off level, however,
grows with $\alpha$ and $p$ and reaches $10^{-4}$ at $\alpha=45^\circ$, $p=8$. On the
folded meshes ($\alpha\ge60^\circ$) the computed solutions for $p\ge5$ are useless
(nodal errors up to $10^{9}$), although the method is exact in exact arithmetic. A
dedicated test (\texttt{results/conditioning.txt}) integrates the \PM\ stiffness
exactly ($n_q\ge p^2+p/2$ points, \cref{sec:quad}) and evaluates the residual
$\bm K\bm u_I-\bm F$ of the exact nodal values, which vanishes in exact arithmetic: it is
$4\cdot10^{-14}$ at $\alpha=0^\circ$, $p=5$, but $2\cdot10^{-2}$ ($\alpha=60^\circ$,
$N=2$, $p=8$) and $49$ ($\alpha=90^\circ$, $N=1$, $p=8$), with
$\sigma_{\min}/\sigma_{\max}(\bm K_{\PM})=10^{-15}$ against $3\cdot10^{-6}$ for \PP.
Distortion immunity therefore holds in exact arithmetic but not in floating point once
the metric interpolation becomes ill-conditioned (Lebesgue constants up to $5\cdot10^{12}$,
\cref{tab:leb}). For $p<5$ the polynomial field behaves like a non-polynomial one
(rates $2,3,4$ for $\alpha\le15^\circ$, \cref{fig:h-poly}).

\begin{table}[t]
  \centering\small
  \caption{Polynomial field (degree 5), \PM: relative nodal error $e_n$ for $p\ge5$
    (exact in exact arithmetic).}
  \label{tab:poly}
  \begin{tabular}{lrrrrrrrr}
\toprule
 & \multicolumn{4}{c}{$N=1$} & \multicolumn{4}{c}{$N=8$}\\
\cmidrule(lr){2-5}\cmidrule(lr){6-9}
 & $p=5$ & $p=6$ & $p=7$ & $p=8$ & $p=5$ & $p=6$ & $p=7$ & $p=8$\\
\midrule
$\alpha=0^\circ$ & 2.3e-14 & 8.7e-14 & 4.0e-12 & 3.3e-11 & 1.4e-13 & 7.7e-14 & 2.8e-13 & 4.6e-13\\
$\alpha=15^\circ$ & 1.1e-11 & 4.9e-13 & 9.6e-12 & 3.7e-11 & 7.4e-14 & 2.2e-13 & 1.6e-11 & 1.6e-11\\
$\alpha=30^\circ$ & 4.4e-12 & 1.7e-12 & 3.5e-11 & 1.8e-9 & 5.0e-11 & 4.3e-11 & 2.5e-9 & 7.7e-9\\
$\alpha=45^\circ$ & 5.3e-12 & 3.2e-9 & 9.8e-9 & 3.1e-6 & 2.1e-8 & 2.0e-7 & 1.7e-5 & 1.1e-4\\
$\alpha=60^\circ$ & 7.9e-11 & 6.1e-9 & 5.9e-8 & 1.5e-6 & 7.4e-1 & 8.5e1 & 3.5e6 & 1.7e9\\
$\alpha=90^\circ$ & 3.4e-7 & 2.2e-5 & 3.5e-3 & 4.0e0 & 5.9e-4 & 3.2e0 & 1.3e2 & 1.9e3\\
\bottomrule
\end{tabular}

\end{table}

\subsection{Non-polynomial fields: convergence}\label{sec:conv}
\Cref{fig:h-smooth,fig:h-lame} show the nodal error, \cref{fig:h-smooth-E} the energy
error for the meshes without folded elements; the rates are collected in
\cref{tab:rates-smooth-E,tab:rates-lame-E,tab:rates-smooth-N,tab:rates-lame-N}. Three
regimes can be distinguished.
\begin{enumerate}
\item $\alpha=0^\circ$ (element maps asymptotically affine): both elements converge
  optimally, $\eE=O(h^p)$ (smooth field: $2.00$, $2.99$, $4.00$, $4.99$, $6.00$ for
  \PM; Lam\'e: $1.97$--$6.62$, approaching $p$ from below as the $r^{-2}$ term is
  resolved), with almost identical errors.
\item $\alpha=15^\circ$: \PM\ keeps the optimal rates for all $p$ (smooth: $2.00$,
  $2.99$, $3.99$, $4.99$, $5.98$; Lam\'e: $1.97$, $2.91$, $3.90$, $4.81$, $5.83$) and
  for $p\le6$ its error exceeds that on the undistorted mesh by only $5$--$12\,\%$
  (smooth) and $5$--$28\,\%$ (Lam\'e) in $\eE$, and by a factor $1.5$--$4.7$ in $e_n$
  (\cref{tab:ratio}). For $p\ge7$ the \PM\ error reaches a round-off floor of about
  $10^{-10}$, two to three orders above the floor on the undistorted mesh. \PP\ loses half of the order: its
  energy rates are $1.08$, $1.23$, $2.21$, $2.03$, $2.94$, $3.39$, $4.05$ (smooth) and,
  for the polynomial field, $0.98$, $0.95$, $1.95$, $1.94$, $2.94$, $2.95$, $3.95$,
  i.e.\ $\lfloor p/2\rfloor$. This is the approximation property of the parametric
  space on maps that are not asymptotically affine: the mid-side node offset is
  $O(h)$, so $\bx(\bxi)$ is quadratic with $O(1)$ relative non-affinity, and
  $Q_p\circ\bx^{-1}$ contains $P_k(\bx)$ only for $2k\le p$ (cf.\
  \citep{Arnold2002Approximation,Ciarlet1972Interpolation}). At $N=8$, $p=6$ the \PM\
  error is $390$ times smaller than the \PP\ error (smooth, $\eE$).
\item $\alpha\ge30^\circ$: \PM\ is optimal only for $p\le3$ at $30^\circ$ (and for $p=2$ at
  $45^\circ$ for the Lam\'e field, rate $1.94$; smooth field $1.11$); for higher $p$ the convergence becomes erratic and even divergent (e.g.\
  smooth, $\alpha=30^\circ$, $p=4$: $\eE=2.2\cdot10^{-4}$ at $N=4$ and $4.3\cdot10^{-4}$ at
  $N=8$; $\alpha=45^\circ$, $p=8$: $\eE=3\cdot10^{-3}$ at $N=8$). On the folded meshes
  ($\alpha\ge60^\circ$) \PM\ fails for $p\ge5$ (\cref{sec:poly}). \PP\ converges
  monotonically at the reduced rate $\approx\lfloor p/2\rfloor$ in $\eE$ (nodal error:
  $\approx\lfloor p/2\rfloor+1$) for all $\alpha\le45^\circ$ and all $p$ (energy
  norm; the nodal error with one exception, $p=5$ at $45^\circ$). On the folded meshes
  $\alpha\ge60^\circ$ its convergence is irregular for $p\le3$, but the nodal error still
  decreases by $1.5$--$3$ orders of magnitude between $N=1$ and $N=8$ for $p\ge4$
  (\cref{fig:h-smooth}).
\end{enumerate}
Under $p$-refinement on a fixed mesh (\cref{fig:p}), \PM\ converges exponentially
for $\alpha\le15^\circ$ down to its round-off floor (reached at $p\approx6$ for $N=8$),
like \PP\ on the undistorted mesh. At $\alpha=30^\circ$ and $45^\circ$ the \PM\ error is
non-monotone in $p$ (smooth field, $N=8$, $\alpha=30^\circ$: $e_n$ jumps from
$2.4\cdot10^{-6}$ at $p=3$ to $2.8\cdot10^{-5}$ at $p=4$ and grows from
$9.4\cdot10^{-11}$ at $p=6$ to $8.0\cdot10^{-9}$ at $p=8$), while the \PP\ error decreases
monotonically, although more slowly than on the undistorted mesh.

\begin{figure}[p]
  \centering
  \begin{tikzpicture}
\begin{groupplot}[group style={group size=6 by 2, horizontal sep=1.5mm, vertical sep=8mm, yticklabels at=edge left, xticklabels at=edge bottom}, conv, width=3.90cm, height=4.6cm, ymin=1e-15, ymax=1e2]
\nextgroupplot[title={\PM, $\alpha=0^\circ$}, ylabel={$e_n$}]
\addplot[pline=2] table[x=h, y=eN2] {\tab/h_smooth_PM_a000.dat};
\addplot[pline=3] table[x=h, y=eN3] {\tab/h_smooth_PM_a000.dat};
\addplot[pline=4] table[x=h, y=eN4] {\tab/h_smooth_PM_a000.dat};
\addplot[pline=5] table[x=h, y=eN5] {\tab/h_smooth_PM_a000.dat};
\addplot[pline=6] table[x=h, y=eN6] {\tab/h_smooth_PM_a000.dat};
\addplot[pline=7] table[x=h, y=eN7] {\tab/h_smooth_PM_a000.dat};
\addplot[pline=8] table[x=h, y=eN8] {\tab/h_smooth_PM_a000.dat};
\nextgroupplot[title={\PM, $\alpha=15^\circ$}]
\addplot[pline=2] table[x=h, y=eN2] {\tab/h_smooth_PM_a015.dat};
\addplot[pline=3] table[x=h, y=eN3] {\tab/h_smooth_PM_a015.dat};
\addplot[pline=4] table[x=h, y=eN4] {\tab/h_smooth_PM_a015.dat};
\addplot[pline=5] table[x=h, y=eN5] {\tab/h_smooth_PM_a015.dat};
\addplot[pline=6] table[x=h, y=eN6] {\tab/h_smooth_PM_a015.dat};
\addplot[pline=7] table[x=h, y=eN7] {\tab/h_smooth_PM_a015.dat};
\addplot[pline=8] table[x=h, y=eN8] {\tab/h_smooth_PM_a015.dat};
\nextgroupplot[title={\PM, $\alpha=30^\circ$}]
\addplot[pline=2] table[x=h, y=eN2] {\tab/h_smooth_PM_a030.dat};
\addplot[pline=3] table[x=h, y=eN3] {\tab/h_smooth_PM_a030.dat};
\addplot[pline=4] table[x=h, y=eN4] {\tab/h_smooth_PM_a030.dat};
\addplot[pline=5] table[x=h, y=eN5] {\tab/h_smooth_PM_a030.dat};
\addplot[pline=6] table[x=h, y=eN6] {\tab/h_smooth_PM_a030.dat};
\addplot[pline=7] table[x=h, y=eN7] {\tab/h_smooth_PM_a030.dat};
\addplot[pline=8] table[x=h, y=eN8] {\tab/h_smooth_PM_a030.dat};
\nextgroupplot[title={\PM, $\alpha=45^\circ$}]
\addplot[pline=2] table[x=h, y=eN2] {\tab/h_smooth_PM_a045.dat};
\addplot[pline=3] table[x=h, y=eN3] {\tab/h_smooth_PM_a045.dat};
\addplot[pline=4] table[x=h, y=eN4] {\tab/h_smooth_PM_a045.dat};
\addplot[pline=5] table[x=h, y=eN5] {\tab/h_smooth_PM_a045.dat};
\addplot[pline=6] table[x=h, y=eN6] {\tab/h_smooth_PM_a045.dat};
\addplot[pline=7] table[x=h, y=eN7] {\tab/h_smooth_PM_a045.dat};
\addplot[pline=8] table[x=h, y=eN8] {\tab/h_smooth_PM_a045.dat};
\nextgroupplot[title={\PM, $\alpha=60^\circ$}]
\addplot[pline=2] table[x=h, y=eN2] {\tab/h_smooth_PM_a060.dat};
\addplot[pline=3] table[x=h, y=eN3] {\tab/h_smooth_PM_a060.dat};
\addplot[pline=4] table[x=h, y=eN4] {\tab/h_smooth_PM_a060.dat};
\addplot[pline=5] table[x=h, y=eN5] {\tab/h_smooth_PM_a060.dat};
\addplot[pline=6] table[x=h, y=eN6] {\tab/h_smooth_PM_a060.dat};
\addplot[pline=7] table[x=h, y=eN7] {\tab/h_smooth_PM_a060.dat};
\addplot[pline=8] table[x=h, y=eN8] {\tab/h_smooth_PM_a060.dat};
\nextgroupplot[title={\PM, $\alpha=90^\circ$}]
\addplot[pline=2] table[x=h, y=eN2] {\tab/h_smooth_PM_a090.dat};
\addplot[pline=3] table[x=h, y=eN3] {\tab/h_smooth_PM_a090.dat};
\addplot[pline=4] table[x=h, y=eN4] {\tab/h_smooth_PM_a090.dat};
\addplot[pline=5] table[x=h, y=eN5] {\tab/h_smooth_PM_a090.dat};
\addplot[pline=6] table[x=h, y=eN6] {\tab/h_smooth_PM_a090.dat};
\addplot[pline=7] table[x=h, y=eN7] {\tab/h_smooth_PM_a090.dat};
\addplot[pline=8] table[x=h, y=eN8] {\tab/h_smooth_PM_a090.dat};
\nextgroupplot[title={\PP, $\alpha=0^\circ$}, ylabel={$e_n$}]
\addplot[pline=2] table[x=h, y=eN2] {\tab/h_smooth_PP_a000.dat};
\addplot[pline=3] table[x=h, y=eN3] {\tab/h_smooth_PP_a000.dat};
\addplot[pline=4] table[x=h, y=eN4] {\tab/h_smooth_PP_a000.dat};
\addplot[pline=5] table[x=h, y=eN5] {\tab/h_smooth_PP_a000.dat};
\addplot[pline=6] table[x=h, y=eN6] {\tab/h_smooth_PP_a000.dat};
\addplot[pline=7] table[x=h, y=eN7] {\tab/h_smooth_PP_a000.dat};
\addplot[pline=8] table[x=h, y=eN8] {\tab/h_smooth_PP_a000.dat};
\nextgroupplot[title={\PP, $\alpha=15^\circ$}]
\addplot[pline=2] table[x=h, y=eN2] {\tab/h_smooth_PP_a015.dat};
\addplot[pline=3] table[x=h, y=eN3] {\tab/h_smooth_PP_a015.dat};
\addplot[pline=4] table[x=h, y=eN4] {\tab/h_smooth_PP_a015.dat};
\addplot[pline=5] table[x=h, y=eN5] {\tab/h_smooth_PP_a015.dat};
\addplot[pline=6] table[x=h, y=eN6] {\tab/h_smooth_PP_a015.dat};
\addplot[pline=7] table[x=h, y=eN7] {\tab/h_smooth_PP_a015.dat};
\addplot[pline=8] table[x=h, y=eN8] {\tab/h_smooth_PP_a015.dat};
\nextgroupplot[title={\PP, $\alpha=30^\circ$}]
\addplot[pline=2] table[x=h, y=eN2] {\tab/h_smooth_PP_a030.dat};
\addplot[pline=3] table[x=h, y=eN3] {\tab/h_smooth_PP_a030.dat};
\addplot[pline=4] table[x=h, y=eN4] {\tab/h_smooth_PP_a030.dat};
\addplot[pline=5] table[x=h, y=eN5] {\tab/h_smooth_PP_a030.dat};
\addplot[pline=6] table[x=h, y=eN6] {\tab/h_smooth_PP_a030.dat};
\addplot[pline=7] table[x=h, y=eN7] {\tab/h_smooth_PP_a030.dat};
\addplot[pline=8] table[x=h, y=eN8] {\tab/h_smooth_PP_a030.dat};
\nextgroupplot[title={\PP, $\alpha=45^\circ$}]
\addplot[pline=2] table[x=h, y=eN2] {\tab/h_smooth_PP_a045.dat};
\addplot[pline=3] table[x=h, y=eN3] {\tab/h_smooth_PP_a045.dat};
\addplot[pline=4] table[x=h, y=eN4] {\tab/h_smooth_PP_a045.dat};
\addplot[pline=5] table[x=h, y=eN5] {\tab/h_smooth_PP_a045.dat};
\addplot[pline=6] table[x=h, y=eN6] {\tab/h_smooth_PP_a045.dat};
\addplot[pline=7] table[x=h, y=eN7] {\tab/h_smooth_PP_a045.dat};
\addplot[pline=8] table[x=h, y=eN8] {\tab/h_smooth_PP_a045.dat};
\nextgroupplot[title={\PP, $\alpha=60^\circ$}]
\addplot[pline=2] table[x=h, y=eN2] {\tab/h_smooth_PP_a060.dat};
\addplot[pline=3] table[x=h, y=eN3] {\tab/h_smooth_PP_a060.dat};
\addplot[pline=4] table[x=h, y=eN4] {\tab/h_smooth_PP_a060.dat};
\addplot[pline=5] table[x=h, y=eN5] {\tab/h_smooth_PP_a060.dat};
\addplot[pline=6] table[x=h, y=eN6] {\tab/h_smooth_PP_a060.dat};
\addplot[pline=7] table[x=h, y=eN7] {\tab/h_smooth_PP_a060.dat};
\addplot[pline=8] table[x=h, y=eN8] {\tab/h_smooth_PP_a060.dat};
\nextgroupplot[title={\PP, $\alpha=90^\circ$}]
\addplot[pline=2] table[x=h, y=eN2] {\tab/h_smooth_PP_a090.dat};
\addplot[pline=3] table[x=h, y=eN3] {\tab/h_smooth_PP_a090.dat};
\addplot[pline=4] table[x=h, y=eN4] {\tab/h_smooth_PP_a090.dat};
\addplot[pline=5] table[x=h, y=eN5] {\tab/h_smooth_PP_a090.dat};
\addplot[pline=6] table[x=h, y=eN6] {\tab/h_smooth_PP_a090.dat};
\addplot[pline=7] table[x=h, y=eN7] {\tab/h_smooth_PP_a090.dat};
\addplot[pline=8] table[x=h, y=eN8] {\tab/h_smooth_PP_a090.dat};
\end{groupplot}
\node[font=\small] at ($(group c1r2.south)!0.5!(group c6r2.south)+(0,-8mm)$) {element size $h\propto 1/N$};
\matrix[draw=gray!50, fill=white, font=\scriptsize, column sep=1mm, inner sep=1.5pt, anchor=south] at ($(group c1r1.north)!0.5!(group c6r1.north)+(0,7mm)$) { \draw[pline=2] (0,0) -- (4mm,0) node[black, right] {$p=2$}; \pgfmatrixnextcell \draw[pline=3] (0,0) -- (4mm,0) node[black, right] {$p=3$}; \pgfmatrixnextcell \draw[pline=4] (0,0) -- (4mm,0) node[black, right] {$p=4$}; \pgfmatrixnextcell \draw[pline=5] (0,0) -- (4mm,0) node[black, right] {$p=5$}; \pgfmatrixnextcell \draw[pline=6] (0,0) -- (4mm,0) node[black, right] {$p=6$}; \pgfmatrixnextcell \draw[pline=7] (0,0) -- (4mm,0) node[black, right] {$p=7$}; \pgfmatrixnextcell \draw[pline=8] (0,0) -- (4mm,0) node[black, right] {$p=8$}; \pgfmatrixnextcell  \\ };

\end{tikzpicture}

  \caption{Smooth manufactured field: nodal error $e_n$ versus element size for all
    distortion angles, \PM\ (top) and \PP\ (bottom).}
  \label{fig:h-smooth}
\end{figure}
\begin{figure}[p]
  \centering
  \begin{tikzpicture}
\begin{groupplot}[group style={group size=6 by 2, horizontal sep=1.5mm, vertical sep=8mm, yticklabels at=edge left, xticklabels at=edge bottom}, conv, width=3.90cm, height=4.6cm, ymin=1e-15, ymax=1e2]
\nextgroupplot[title={\PM, $\alpha=0^\circ$}, ylabel={$e_n$}]
\addplot[pline=2] table[x=h, y=eN2] {\tab/h_lame_PM_a000.dat};
\addplot[pline=3] table[x=h, y=eN3] {\tab/h_lame_PM_a000.dat};
\addplot[pline=4] table[x=h, y=eN4] {\tab/h_lame_PM_a000.dat};
\addplot[pline=5] table[x=h, y=eN5] {\tab/h_lame_PM_a000.dat};
\addplot[pline=6] table[x=h, y=eN6] {\tab/h_lame_PM_a000.dat};
\addplot[pline=7] table[x=h, y=eN7] {\tab/h_lame_PM_a000.dat};
\addplot[pline=8] table[x=h, y=eN8] {\tab/h_lame_PM_a000.dat};
\nextgroupplot[title={\PM, $\alpha=15^\circ$}]
\addplot[pline=2] table[x=h, y=eN2] {\tab/h_lame_PM_a015.dat};
\addplot[pline=3] table[x=h, y=eN3] {\tab/h_lame_PM_a015.dat};
\addplot[pline=4] table[x=h, y=eN4] {\tab/h_lame_PM_a015.dat};
\addplot[pline=5] table[x=h, y=eN5] {\tab/h_lame_PM_a015.dat};
\addplot[pline=6] table[x=h, y=eN6] {\tab/h_lame_PM_a015.dat};
\addplot[pline=7] table[x=h, y=eN7] {\tab/h_lame_PM_a015.dat};
\addplot[pline=8] table[x=h, y=eN8] {\tab/h_lame_PM_a015.dat};
\nextgroupplot[title={\PM, $\alpha=30^\circ$}]
\addplot[pline=2] table[x=h, y=eN2] {\tab/h_lame_PM_a030.dat};
\addplot[pline=3] table[x=h, y=eN3] {\tab/h_lame_PM_a030.dat};
\addplot[pline=4] table[x=h, y=eN4] {\tab/h_lame_PM_a030.dat};
\addplot[pline=5] table[x=h, y=eN5] {\tab/h_lame_PM_a030.dat};
\addplot[pline=6] table[x=h, y=eN6] {\tab/h_lame_PM_a030.dat};
\addplot[pline=7] table[x=h, y=eN7] {\tab/h_lame_PM_a030.dat};
\addplot[pline=8] table[x=h, y=eN8] {\tab/h_lame_PM_a030.dat};
\nextgroupplot[title={\PM, $\alpha=45^\circ$}]
\addplot[pline=2] table[x=h, y=eN2] {\tab/h_lame_PM_a045.dat};
\addplot[pline=3] table[x=h, y=eN3] {\tab/h_lame_PM_a045.dat};
\addplot[pline=4] table[x=h, y=eN4] {\tab/h_lame_PM_a045.dat};
\addplot[pline=5] table[x=h, y=eN5] {\tab/h_lame_PM_a045.dat};
\addplot[pline=6] table[x=h, y=eN6] {\tab/h_lame_PM_a045.dat};
\addplot[pline=7] table[x=h, y=eN7] {\tab/h_lame_PM_a045.dat};
\addplot[pline=8] table[x=h, y=eN8] {\tab/h_lame_PM_a045.dat};
\nextgroupplot[title={\PM, $\alpha=60^\circ$}]
\addplot[pline=2] table[x=h, y=eN2] {\tab/h_lame_PM_a060.dat};
\addplot[pline=3] table[x=h, y=eN3] {\tab/h_lame_PM_a060.dat};
\addplot[pline=4] table[x=h, y=eN4] {\tab/h_lame_PM_a060.dat};
\addplot[pline=5] table[x=h, y=eN5] {\tab/h_lame_PM_a060.dat};
\addplot[pline=6] table[x=h, y=eN6] {\tab/h_lame_PM_a060.dat};
\addplot[pline=7] table[x=h, y=eN7] {\tab/h_lame_PM_a060.dat};
\addplot[pline=8] table[x=h, y=eN8] {\tab/h_lame_PM_a060.dat};
\nextgroupplot[title={\PM, $\alpha=90^\circ$}]
\addplot[pline=2] table[x=h, y=eN2] {\tab/h_lame_PM_a090.dat};
\addplot[pline=3] table[x=h, y=eN3] {\tab/h_lame_PM_a090.dat};
\addplot[pline=4] table[x=h, y=eN4] {\tab/h_lame_PM_a090.dat};
\addplot[pline=5] table[x=h, y=eN5] {\tab/h_lame_PM_a090.dat};
\addplot[pline=6] table[x=h, y=eN6] {\tab/h_lame_PM_a090.dat};
\addplot[pline=7] table[x=h, y=eN7] {\tab/h_lame_PM_a090.dat};
\addplot[pline=8] table[x=h, y=eN8] {\tab/h_lame_PM_a090.dat};
\nextgroupplot[title={\PP, $\alpha=0^\circ$}, ylabel={$e_n$}]
\addplot[pline=2] table[x=h, y=eN2] {\tab/h_lame_PP_a000.dat};
\addplot[pline=3] table[x=h, y=eN3] {\tab/h_lame_PP_a000.dat};
\addplot[pline=4] table[x=h, y=eN4] {\tab/h_lame_PP_a000.dat};
\addplot[pline=5] table[x=h, y=eN5] {\tab/h_lame_PP_a000.dat};
\addplot[pline=6] table[x=h, y=eN6] {\tab/h_lame_PP_a000.dat};
\addplot[pline=7] table[x=h, y=eN7] {\tab/h_lame_PP_a000.dat};
\addplot[pline=8] table[x=h, y=eN8] {\tab/h_lame_PP_a000.dat};
\nextgroupplot[title={\PP, $\alpha=15^\circ$}]
\addplot[pline=2] table[x=h, y=eN2] {\tab/h_lame_PP_a015.dat};
\addplot[pline=3] table[x=h, y=eN3] {\tab/h_lame_PP_a015.dat};
\addplot[pline=4] table[x=h, y=eN4] {\tab/h_lame_PP_a015.dat};
\addplot[pline=5] table[x=h, y=eN5] {\tab/h_lame_PP_a015.dat};
\addplot[pline=6] table[x=h, y=eN6] {\tab/h_lame_PP_a015.dat};
\addplot[pline=7] table[x=h, y=eN7] {\tab/h_lame_PP_a015.dat};
\addplot[pline=8] table[x=h, y=eN8] {\tab/h_lame_PP_a015.dat};
\nextgroupplot[title={\PP, $\alpha=30^\circ$}]
\addplot[pline=2] table[x=h, y=eN2] {\tab/h_lame_PP_a030.dat};
\addplot[pline=3] table[x=h, y=eN3] {\tab/h_lame_PP_a030.dat};
\addplot[pline=4] table[x=h, y=eN4] {\tab/h_lame_PP_a030.dat};
\addplot[pline=5] table[x=h, y=eN5] {\tab/h_lame_PP_a030.dat};
\addplot[pline=6] table[x=h, y=eN6] {\tab/h_lame_PP_a030.dat};
\addplot[pline=7] table[x=h, y=eN7] {\tab/h_lame_PP_a030.dat};
\addplot[pline=8] table[x=h, y=eN8] {\tab/h_lame_PP_a030.dat};
\nextgroupplot[title={\PP, $\alpha=45^\circ$}]
\addplot[pline=2] table[x=h, y=eN2] {\tab/h_lame_PP_a045.dat};
\addplot[pline=3] table[x=h, y=eN3] {\tab/h_lame_PP_a045.dat};
\addplot[pline=4] table[x=h, y=eN4] {\tab/h_lame_PP_a045.dat};
\addplot[pline=5] table[x=h, y=eN5] {\tab/h_lame_PP_a045.dat};
\addplot[pline=6] table[x=h, y=eN6] {\tab/h_lame_PP_a045.dat};
\addplot[pline=7] table[x=h, y=eN7] {\tab/h_lame_PP_a045.dat};
\addplot[pline=8] table[x=h, y=eN8] {\tab/h_lame_PP_a045.dat};
\nextgroupplot[title={\PP, $\alpha=60^\circ$}]
\addplot[pline=2] table[x=h, y=eN2] {\tab/h_lame_PP_a060.dat};
\addplot[pline=3] table[x=h, y=eN3] {\tab/h_lame_PP_a060.dat};
\addplot[pline=4] table[x=h, y=eN4] {\tab/h_lame_PP_a060.dat};
\addplot[pline=5] table[x=h, y=eN5] {\tab/h_lame_PP_a060.dat};
\addplot[pline=6] table[x=h, y=eN6] {\tab/h_lame_PP_a060.dat};
\addplot[pline=7] table[x=h, y=eN7] {\tab/h_lame_PP_a060.dat};
\addplot[pline=8] table[x=h, y=eN8] {\tab/h_lame_PP_a060.dat};
\nextgroupplot[title={\PP, $\alpha=90^\circ$}]
\addplot[pline=2] table[x=h, y=eN2] {\tab/h_lame_PP_a090.dat};
\addplot[pline=3] table[x=h, y=eN3] {\tab/h_lame_PP_a090.dat};
\addplot[pline=4] table[x=h, y=eN4] {\tab/h_lame_PP_a090.dat};
\addplot[pline=5] table[x=h, y=eN5] {\tab/h_lame_PP_a090.dat};
\addplot[pline=6] table[x=h, y=eN6] {\tab/h_lame_PP_a090.dat};
\addplot[pline=7] table[x=h, y=eN7] {\tab/h_lame_PP_a090.dat};
\addplot[pline=8] table[x=h, y=eN8] {\tab/h_lame_PP_a090.dat};
\end{groupplot}
\node[font=\small] at ($(group c1r2.south)!0.5!(group c6r2.south)+(0,-8mm)$) {element size $h\propto 1/N$};
\matrix[draw=gray!50, fill=white, font=\scriptsize, column sep=1mm, inner sep=1.5pt, anchor=south] at ($(group c1r1.north)!0.5!(group c6r1.north)+(0,7mm)$) { \draw[pline=2] (0,0) -- (4mm,0) node[black, right] {$p=2$}; \pgfmatrixnextcell \draw[pline=3] (0,0) -- (4mm,0) node[black, right] {$p=3$}; \pgfmatrixnextcell \draw[pline=4] (0,0) -- (4mm,0) node[black, right] {$p=4$}; \pgfmatrixnextcell \draw[pline=5] (0,0) -- (4mm,0) node[black, right] {$p=5$}; \pgfmatrixnextcell \draw[pline=6] (0,0) -- (4mm,0) node[black, right] {$p=6$}; \pgfmatrixnextcell \draw[pline=7] (0,0) -- (4mm,0) node[black, right] {$p=7$}; \pgfmatrixnextcell \draw[pline=8] (0,0) -- (4mm,0) node[black, right] {$p=8$}; \pgfmatrixnextcell  \\ };

\end{tikzpicture}

  \caption{Thick cylinder (Lam\'e): nodal error $e_n$ versus element size, \PM\ (top)
    and \PP\ (bottom).}
  \label{fig:h-lame}
\end{figure}
\begin{figure}[t]
  \centering
  \begin{tikzpicture}
\begin{groupplot}[group style={group size=6 by 3, horizontal sep=1.5mm, vertical sep=8mm, yticklabels at=edge left, xticklabels at=edge bottom}, conv, width=3.90cm, height=4.4cm, ymin=1e-15, ymax=1e2]
\nextgroupplot[title={MMS $\alpha=0^\circ$}, ylabel={\PM: $\eE$}]
\addplot[pline=2] table[x=h, y=eE2] {\tab/h_smooth_PM_a000.dat};
\addplot[pline=3] table[x=h, y=eE3] {\tab/h_smooth_PM_a000.dat};
\addplot[pline=4] table[x=h, y=eE4] {\tab/h_smooth_PM_a000.dat};
\addplot[pline=5] table[x=h, y=eE5] {\tab/h_smooth_PM_a000.dat};
\addplot[pline=6] table[x=h, y=eE6] {\tab/h_smooth_PM_a000.dat};
\addplot[pline=7] table[x=h, y=eE7] {\tab/h_smooth_PM_a000.dat};
\addplot[pline=8] table[x=h, y=eE8] {\tab/h_smooth_PM_a000.dat};
\nextgroupplot[title={MMS $\alpha=15^\circ$}]
\addplot[pline=2] table[x=h, y=eE2] {\tab/h_smooth_PM_a015.dat};
\addplot[pline=3] table[x=h, y=eE3] {\tab/h_smooth_PM_a015.dat};
\addplot[pline=4] table[x=h, y=eE4] {\tab/h_smooth_PM_a015.dat};
\addplot[pline=5] table[x=h, y=eE5] {\tab/h_smooth_PM_a015.dat};
\addplot[pline=6] table[x=h, y=eE6] {\tab/h_smooth_PM_a015.dat};
\addplot[pline=7] table[x=h, y=eE7] {\tab/h_smooth_PM_a015.dat};
\addplot[pline=8] table[x=h, y=eE8] {\tab/h_smooth_PM_a015.dat};
\nextgroupplot[title={MMS $\alpha=30^\circ$}]
\addplot[pline=2] table[x=h, y=eE2] {\tab/h_smooth_PM_a030.dat};
\addplot[pline=3] table[x=h, y=eE3] {\tab/h_smooth_PM_a030.dat};
\addplot[pline=4] table[x=h, y=eE4] {\tab/h_smooth_PM_a030.dat};
\addplot[pline=5] table[x=h, y=eE5] {\tab/h_smooth_PM_a030.dat};
\addplot[pline=6] table[x=h, y=eE6] {\tab/h_smooth_PM_a030.dat};
\addplot[pline=7] table[x=h, y=eE7] {\tab/h_smooth_PM_a030.dat};
\addplot[pline=8] table[x=h, y=eE8] {\tab/h_smooth_PM_a030.dat};
\nextgroupplot[title={Lam\'e $\alpha=0^\circ$}]
\addplot[pline=2] table[x=h, y=eE2] {\tab/h_lame_PM_a000.dat};
\addplot[pline=3] table[x=h, y=eE3] {\tab/h_lame_PM_a000.dat};
\addplot[pline=4] table[x=h, y=eE4] {\tab/h_lame_PM_a000.dat};
\addplot[pline=5] table[x=h, y=eE5] {\tab/h_lame_PM_a000.dat};
\addplot[pline=6] table[x=h, y=eE6] {\tab/h_lame_PM_a000.dat};
\addplot[pline=7] table[x=h, y=eE7] {\tab/h_lame_PM_a000.dat};
\addplot[pline=8] table[x=h, y=eE8] {\tab/h_lame_PM_a000.dat};
\nextgroupplot[title={Lam\'e $\alpha=15^\circ$}]
\addplot[pline=2] table[x=h, y=eE2] {\tab/h_lame_PM_a015.dat};
\addplot[pline=3] table[x=h, y=eE3] {\tab/h_lame_PM_a015.dat};
\addplot[pline=4] table[x=h, y=eE4] {\tab/h_lame_PM_a015.dat};
\addplot[pline=5] table[x=h, y=eE5] {\tab/h_lame_PM_a015.dat};
\addplot[pline=6] table[x=h, y=eE6] {\tab/h_lame_PM_a015.dat};
\addplot[pline=7] table[x=h, y=eE7] {\tab/h_lame_PM_a015.dat};
\addplot[pline=8] table[x=h, y=eE8] {\tab/h_lame_PM_a015.dat};
\nextgroupplot[title={Lam\'e $\alpha=30^\circ$}]
\addplot[pline=2] table[x=h, y=eE2] {\tab/h_lame_PM_a030.dat};
\addplot[pline=3] table[x=h, y=eE3] {\tab/h_lame_PM_a030.dat};
\addplot[pline=4] table[x=h, y=eE4] {\tab/h_lame_PM_a030.dat};
\addplot[pline=5] table[x=h, y=eE5] {\tab/h_lame_PM_a030.dat};
\addplot[pline=6] table[x=h, y=eE6] {\tab/h_lame_PM_a030.dat};
\addplot[pline=7] table[x=h, y=eE7] {\tab/h_lame_PM_a030.dat};
\addplot[pline=8] table[x=h, y=eE8] {\tab/h_lame_PM_a030.dat};
\nextgroupplot[title={MMS $\alpha=0^\circ$}, ylabel={\PH: $\eE$}]
\addplot[pline=2] table[x=h, y=eE2] {\tab/h_smooth_PH_a000.dat};
\addplot[pline=3] table[x=h, y=eE3] {\tab/h_smooth_PH_a000.dat};
\addplot[pline=4] table[x=h, y=eE4] {\tab/h_smooth_PH_a000.dat};
\addplot[pline=5] table[x=h, y=eE5] {\tab/h_smooth_PH_a000.dat};
\addplot[pline=6] table[x=h, y=eE6] {\tab/h_smooth_PH_a000.dat};
\addplot[pline=7] table[x=h, y=eE7] {\tab/h_smooth_PH_a000.dat};
\addplot[pline=8] table[x=h, y=eE8] {\tab/h_smooth_PH_a000.dat};
\nextgroupplot[title={MMS $\alpha=15^\circ$}]
\addplot[pline=2] table[x=h, y=eE2] {\tab/h_smooth_PH_a015.dat};
\addplot[pline=3] table[x=h, y=eE3] {\tab/h_smooth_PH_a015.dat};
\addplot[pline=4] table[x=h, y=eE4] {\tab/h_smooth_PH_a015.dat};
\addplot[pline=5] table[x=h, y=eE5] {\tab/h_smooth_PH_a015.dat};
\addplot[pline=6] table[x=h, y=eE6] {\tab/h_smooth_PH_a015.dat};
\addplot[pline=7] table[x=h, y=eE7] {\tab/h_smooth_PH_a015.dat};
\addplot[pline=8] table[x=h, y=eE8] {\tab/h_smooth_PH_a015.dat};
\nextgroupplot[title={MMS $\alpha=30^\circ$}]
\addplot[pline=2] table[x=h, y=eE2] {\tab/h_smooth_PH_a030.dat};
\addplot[pline=3] table[x=h, y=eE3] {\tab/h_smooth_PH_a030.dat};
\addplot[pline=4] table[x=h, y=eE4] {\tab/h_smooth_PH_a030.dat};
\addplot[pline=5] table[x=h, y=eE5] {\tab/h_smooth_PH_a030.dat};
\addplot[pline=6] table[x=h, y=eE6] {\tab/h_smooth_PH_a030.dat};
\addplot[pline=7] table[x=h, y=eE7] {\tab/h_smooth_PH_a030.dat};
\addplot[pline=8] table[x=h, y=eE8] {\tab/h_smooth_PH_a030.dat};
\nextgroupplot[title={Lam\'e $\alpha=0^\circ$}]
\addplot[pline=2] table[x=h, y=eE2] {\tab/h_lame_PH_a000.dat};
\addplot[pline=3] table[x=h, y=eE3] {\tab/h_lame_PH_a000.dat};
\addplot[pline=4] table[x=h, y=eE4] {\tab/h_lame_PH_a000.dat};
\addplot[pline=5] table[x=h, y=eE5] {\tab/h_lame_PH_a000.dat};
\addplot[pline=6] table[x=h, y=eE6] {\tab/h_lame_PH_a000.dat};
\addplot[pline=7] table[x=h, y=eE7] {\tab/h_lame_PH_a000.dat};
\addplot[pline=8] table[x=h, y=eE8] {\tab/h_lame_PH_a000.dat};
\nextgroupplot[title={Lam\'e $\alpha=15^\circ$}]
\addplot[pline=2] table[x=h, y=eE2] {\tab/h_lame_PH_a015.dat};
\addplot[pline=3] table[x=h, y=eE3] {\tab/h_lame_PH_a015.dat};
\addplot[pline=4] table[x=h, y=eE4] {\tab/h_lame_PH_a015.dat};
\addplot[pline=5] table[x=h, y=eE5] {\tab/h_lame_PH_a015.dat};
\addplot[pline=6] table[x=h, y=eE6] {\tab/h_lame_PH_a015.dat};
\addplot[pline=7] table[x=h, y=eE7] {\tab/h_lame_PH_a015.dat};
\addplot[pline=8] table[x=h, y=eE8] {\tab/h_lame_PH_a015.dat};
\nextgroupplot[title={Lam\'e $\alpha=30^\circ$}]
\addplot[pline=2] table[x=h, y=eE2] {\tab/h_lame_PH_a030.dat};
\addplot[pline=3] table[x=h, y=eE3] {\tab/h_lame_PH_a030.dat};
\addplot[pline=4] table[x=h, y=eE4] {\tab/h_lame_PH_a030.dat};
\addplot[pline=5] table[x=h, y=eE5] {\tab/h_lame_PH_a030.dat};
\addplot[pline=6] table[x=h, y=eE6] {\tab/h_lame_PH_a030.dat};
\addplot[pline=7] table[x=h, y=eE7] {\tab/h_lame_PH_a030.dat};
\addplot[pline=8] table[x=h, y=eE8] {\tab/h_lame_PH_a030.dat};
\nextgroupplot[title={MMS $\alpha=0^\circ$}, ylabel={\PP: $\eE$}]
\addplot[pline=2] table[x=h, y=eE2] {\tab/h_smooth_PP_a000.dat};
\addplot[pline=3] table[x=h, y=eE3] {\tab/h_smooth_PP_a000.dat};
\addplot[pline=4] table[x=h, y=eE4] {\tab/h_smooth_PP_a000.dat};
\addplot[pline=5] table[x=h, y=eE5] {\tab/h_smooth_PP_a000.dat};
\addplot[pline=6] table[x=h, y=eE6] {\tab/h_smooth_PP_a000.dat};
\addplot[pline=7] table[x=h, y=eE7] {\tab/h_smooth_PP_a000.dat};
\addplot[pline=8] table[x=h, y=eE8] {\tab/h_smooth_PP_a000.dat};
\nextgroupplot[title={MMS $\alpha=15^\circ$}]
\addplot[pline=2] table[x=h, y=eE2] {\tab/h_smooth_PP_a015.dat};
\addplot[pline=3] table[x=h, y=eE3] {\tab/h_smooth_PP_a015.dat};
\addplot[pline=4] table[x=h, y=eE4] {\tab/h_smooth_PP_a015.dat};
\addplot[pline=5] table[x=h, y=eE5] {\tab/h_smooth_PP_a015.dat};
\addplot[pline=6] table[x=h, y=eE6] {\tab/h_smooth_PP_a015.dat};
\addplot[pline=7] table[x=h, y=eE7] {\tab/h_smooth_PP_a015.dat};
\addplot[pline=8] table[x=h, y=eE8] {\tab/h_smooth_PP_a015.dat};
\nextgroupplot[title={MMS $\alpha=30^\circ$}]
\addplot[pline=2] table[x=h, y=eE2] {\tab/h_smooth_PP_a030.dat};
\addplot[pline=3] table[x=h, y=eE3] {\tab/h_smooth_PP_a030.dat};
\addplot[pline=4] table[x=h, y=eE4] {\tab/h_smooth_PP_a030.dat};
\addplot[pline=5] table[x=h, y=eE5] {\tab/h_smooth_PP_a030.dat};
\addplot[pline=6] table[x=h, y=eE6] {\tab/h_smooth_PP_a030.dat};
\addplot[pline=7] table[x=h, y=eE7] {\tab/h_smooth_PP_a030.dat};
\addplot[pline=8] table[x=h, y=eE8] {\tab/h_smooth_PP_a030.dat};
\nextgroupplot[title={Lam\'e $\alpha=0^\circ$}]
\addplot[pline=2] table[x=h, y=eE2] {\tab/h_lame_PP_a000.dat};
\addplot[pline=3] table[x=h, y=eE3] {\tab/h_lame_PP_a000.dat};
\addplot[pline=4] table[x=h, y=eE4] {\tab/h_lame_PP_a000.dat};
\addplot[pline=5] table[x=h, y=eE5] {\tab/h_lame_PP_a000.dat};
\addplot[pline=6] table[x=h, y=eE6] {\tab/h_lame_PP_a000.dat};
\addplot[pline=7] table[x=h, y=eE7] {\tab/h_lame_PP_a000.dat};
\addplot[pline=8] table[x=h, y=eE8] {\tab/h_lame_PP_a000.dat};
\nextgroupplot[title={Lam\'e $\alpha=15^\circ$}]
\addplot[pline=2] table[x=h, y=eE2] {\tab/h_lame_PP_a015.dat};
\addplot[pline=3] table[x=h, y=eE3] {\tab/h_lame_PP_a015.dat};
\addplot[pline=4] table[x=h, y=eE4] {\tab/h_lame_PP_a015.dat};
\addplot[pline=5] table[x=h, y=eE5] {\tab/h_lame_PP_a015.dat};
\addplot[pline=6] table[x=h, y=eE6] {\tab/h_lame_PP_a015.dat};
\addplot[pline=7] table[x=h, y=eE7] {\tab/h_lame_PP_a015.dat};
\addplot[pline=8] table[x=h, y=eE8] {\tab/h_lame_PP_a015.dat};
\nextgroupplot[title={Lam\'e $\alpha=30^\circ$}]
\addplot[pline=2] table[x=h, y=eE2] {\tab/h_lame_PP_a030.dat};
\addplot[pline=3] table[x=h, y=eE3] {\tab/h_lame_PP_a030.dat};
\addplot[pline=4] table[x=h, y=eE4] {\tab/h_lame_PP_a030.dat};
\addplot[pline=5] table[x=h, y=eE5] {\tab/h_lame_PP_a030.dat};
\addplot[pline=6] table[x=h, y=eE6] {\tab/h_lame_PP_a030.dat};
\addplot[pline=7] table[x=h, y=eE7] {\tab/h_lame_PP_a030.dat};
\addplot[pline=8] table[x=h, y=eE8] {\tab/h_lame_PP_a030.dat};
\end{groupplot}
\node[font=\small] at ($(group c1r3.south)!0.5!(group c6r3.south)+(0,-8mm)$) {element size $h\propto 1/N$};
\matrix[draw=gray!50, fill=white, font=\scriptsize, column sep=1mm, inner sep=1.5pt, anchor=south] at ($(group c1r1.north)!0.5!(group c6r1.north)+(0,7mm)$) { \draw[pline=2] (0,0) -- (4mm,0) node[black, right] {$p=2$}; \pgfmatrixnextcell \draw[pline=3] (0,0) -- (4mm,0) node[black, right] {$p=3$}; \pgfmatrixnextcell \draw[pline=4] (0,0) -- (4mm,0) node[black, right] {$p=4$}; \pgfmatrixnextcell \draw[pline=5] (0,0) -- (4mm,0) node[black, right] {$p=5$}; \pgfmatrixnextcell \draw[pline=6] (0,0) -- (4mm,0) node[black, right] {$p=6$}; \pgfmatrixnextcell \draw[pline=7] (0,0) -- (4mm,0) node[black, right] {$p=7$}; \pgfmatrixnextcell \draw[pline=8] (0,0) -- (4mm,0) node[black, right] {$p=8$}; \pgfmatrixnextcell  \\ };

\end{tikzpicture}

  \caption{Energy error $\eE$ on the meshes without folded elements
    ($\alpha\le30^\circ$): smooth field (left three columns) and Lam\'e (right three
    columns); \PM\ top, \PP\ bottom.}
  \label{fig:h-smooth-E}
\end{figure}

\begin{table}[p]
  \centering\small
  \begin{minipage}[t]{0.48\linewidth}\centering
    \captionof{table}{Rates of $\eE$, smooth field ($N=4\to8$).}\label{tab:rates-smooth-E}
    \resizebox{\linewidth}{!}{\begin{tabular}{lrrrrrrr}
\toprule
& \multicolumn{7}{c}{$p$}\\
\cmidrule(lr){2-8}
 & 2 & 3 & 4 & 5 & 6 & 7 & 8\\
\midrule
\multicolumn{8}{l}{\PM\ (unsymmetric, skew coordinates)}\\
$\alpha=0^\circ$ & 2.00 & 2.99 & 4.00 & 4.99 & 6.00 & -- & --\\
$\alpha=15^\circ$ & 2.00 & 2.99 & 3.99 & 4.99 & 5.98 & -- & --\\
$\alpha=30^\circ$ & 2.00 & 2.86 & -0.98 & 2.41 & -- & -- & --\\
\midrule
\multicolumn{8}{l}{\PH\ (unsymmetric, hybrid space)}\\
$\alpha=0^\circ$ & 2.00 & 2.99 & 4.00 & 4.99 & 6.00 & -- & --\\
$\alpha=15^\circ$ & 2.00 & 3.00 & 4.00 & 4.99 & 6.01 & -- & --\\
$\alpha=30^\circ$ & 2.00 & 3.00 & 4.01 & 5.00 & 6.02 & -- & --\\
\midrule
\multicolumn{8}{l}{\PP\ (conventional)}\\
$\alpha=0^\circ$ & 2.02 & 3.19 & 4.22 & 5.37 & 6.39 & 7.43 & --\\
$\alpha=15^\circ$ & 1.08 & 1.23 & 2.21 & 2.03 & 2.94 & 3.39 & 4.05\\
$\alpha=30^\circ$ & 1.02 & 1.06 & 2.01 & 1.98 & 2.98 & 3.11 & 3.94\\
\bottomrule
\end{tabular}
}
  \end{minipage}\hfill
  \begin{minipage}[t]{0.48\linewidth}\centering
    \captionof{table}{Rates of $\eE$, Lam\'e ($N=4\to8$).}\label{tab:rates-lame-E}
    \resizebox{\linewidth}{!}{\begin{tabular}{lrrrrrrr}
\toprule
& \multicolumn{7}{c}{$p$}\\
\cmidrule(lr){2-8}
 & 2 & 3 & 4 & 5 & 6 & 7 & 8\\
\midrule
\multicolumn{8}{l}{\PM\ (unsymmetric, skew coordinates)}\\
$\alpha=0^\circ$ & 1.97 & 2.90 & 3.89 & 4.78 & 5.77 & 6.62 & --\\
$\alpha=15^\circ$ & 1.97 & 2.91 & 3.90 & 4.81 & 5.83 & -- & --\\
$\alpha=30^\circ$ & 1.97 & 2.94 & 0.76 & 2.99 & 4.56 & -- & --\\
\midrule
\multicolumn{8}{l}{\PH\ (unsymmetric, hybrid space)}\\
$\alpha=0^\circ$ & 1.96 & 2.91 & 3.89 & 4.81 & 5.79 & 6.69 & --\\
$\alpha=15^\circ$ & 1.96 & 2.94 & 3.94 & 4.92 & 5.91 & 6.88 & 7.87\\
$\alpha=30^\circ$ & 1.96 & 2.96 & 3.94 & 4.94 & 5.92 & 6.90 & 7.87\\
\midrule
\multicolumn{8}{l}{\PP\ (conventional)}\\
$\alpha=0^\circ$ & 1.91 & 2.89 & 3.85 & 4.79 & 5.75 & 6.67 & 7.59\\
$\alpha=15^\circ$ & 1.07 & 1.43 & 2.16 & 2.81 & 3.49 & 4.19 & 4.78\\
$\alpha=30^\circ$ & 1.03 & 1.21 & 2.07 & 2.51 & 3.20 & 3.85 & 4.42\\
\bottomrule
\end{tabular}
}
  \end{minipage}

  \vspace{4mm}
  \begin{minipage}[t]{0.48\linewidth}\centering
    \captionof{table}{Rates of $e_n$, smooth field ($N=4\to8$).}\label{tab:rates-smooth-N}
    \resizebox{\linewidth}{!}{\begin{tabular}{lrrrrrrr}
\toprule
& \multicolumn{7}{c}{$p$}\\
\cmidrule(lr){2-8}
 & 2 & 3 & 4 & 5 & 6 & 7 & 8\\
\midrule
\multicolumn{8}{l}{\PM\ (unsymmetric)}\\
$\alpha=0^\circ$ & 2.70 & 3.80 & 4.86 & 5.84 & 6.92 & -- & --\\
$\alpha=15^\circ$ & 2.86 & 3.83 & 4.81 & 5.79 & 6.81 & -- & --\\
$\alpha=30^\circ$ & 2.52 & 3.30 & -1.00 & 1.97 & -- & -- & --\\
$\alpha=45^\circ$$^\dagger$ & 2.31 & 1.12 & 1.62 & -1.00 & -- & -- & --\\
$\alpha=60^\circ$$^\dagger$ & -0.37 & 3.52 & -2.08 & -- & -- & -- & --\\
$\alpha=90^\circ$$^\dagger$ & 1.20 & -1.23 & 4.76 & -- & -- & -- & --\\
\midrule
\multicolumn{8}{l}{\PP\ (conventional)}\\
$\alpha=0^\circ$ & 2.91 & 3.80 & 4.73 & 5.78 & 6.79 & -- & --\\
$\alpha=15^\circ$ & 1.87 & 2.41 & 3.14 & 2.80 & 3.66 & 4.13 & 4.49\\
$\alpha=30^\circ$ & 1.87 & 2.00 & 2.65 & 2.65 & 3.61 & 3.79 & 4.59\\
$\alpha=45^\circ$$^\dagger$ & 1.65 & 1.62 & 2.28 & 5.03 & 3.49 & 2.84 & 3.73\\
$\alpha=60^\circ$$^\dagger$ & 5.79 & 6.91 & 0.42 & 1.39 & 3.57 & 3.10 & 4.01\\
$\alpha=90^\circ$$^\dagger$ & -0.26 & -2.79 & 2.45 & 1.51 & 3.46 & 1.62 & 7.08\\
\bottomrule
\end{tabular}
}
  \end{minipage}\hfill
  \begin{minipage}[t]{0.48\linewidth}\centering
    \captionof{table}{Rates of $e_n$, Lam\'e ($N=4\to8$).}\label{tab:rates-lame-N}
    \resizebox{\linewidth}{!}{\begin{tabular}{lrrrrrrr}
\toprule
& \multicolumn{7}{c}{$p$}\\
\cmidrule(lr){2-8}
 & 2 & 3 & 4 & 5 & 6 & 7 & 8\\
\midrule
\multicolumn{8}{l}{\PM\ (unsymmetric, skew coordinates)}\\
$\alpha=0^\circ$ & 3.48 & 3.98 & 5.21 & 5.91 & 7.21 & -- & --\\
$\alpha=15^\circ$ & 3.03 & 3.90 & 4.96 & 5.89 & 6.59 & -- & --\\
$\alpha=30^\circ$ & 3.19 & 3.92 & 0.51 & 2.85 & 5.08 & -- & --\\
$\alpha=45^\circ$$^\dagger$ & 2.57 & 1.95 & 0.87 & 3.86 & 2.93 & -1.81 & --\\
$\alpha=60^\circ$$^\dagger$ & 4.35 & 2.84 & 6.22 & 2.08 & -- & -- & --\\
$\alpha=90^\circ$$^\dagger$ & 2.23 & 1.42 & 5.90 & 2.08 & 6.06 & -- & --\\
\midrule
\multicolumn{8}{l}{\PH\ (unsymmetric, hybrid space)}\\
$\alpha=0^\circ$ & 3.19 & 4.01 & 5.02 & 5.52 & 7.04 & -- & --\\
$\alpha=15^\circ$ & 2.12 & 3.79 & 4.78 & 5.50 & 6.44 & 7.35 & --\\
$\alpha=30^\circ$ & 2.10 & 3.75 & 4.64 & 5.54 & 6.52 & 7.56 & 7.79\\
$\alpha=45^\circ$$^\dagger$ & 2.09 & 3.72 & 4.48 & 5.74 & 6.52 & 7.32 & --\\
$\alpha=60^\circ$$^\dagger$ & -1.73 & 3.63 & 4.35 & 5.96 & 6.25 & -- & --\\
$\alpha=90^\circ$$^\dagger$ & -4.20 & 1.30 & 1.15 & 4.94 & 10.87 & -- & --\\
\midrule
\multicolumn{8}{l}{\PP\ (conventional)}\\
$\alpha=0^\circ$ & 3.47 & 4.06 & 5.31 & 6.49 & 6.92 & 3.27 & --\\
$\alpha=15^\circ$ & 2.02 & 2.70 & 2.96 & 3.15 & 3.90 & 4.58 & 5.13\\
$\alpha=30^\circ$ & 1.88 & 2.36 & 3.12 & 2.91 & 3.69 & 4.46 & 5.11\\
$\alpha=45^\circ$$^\dagger$ & 1.81 & 2.17 & 3.00 & 2.78 & 3.59 & 4.13 & 4.67\\
$\alpha=60^\circ$$^\dagger$ & 2.37 & -0.59 & 1.74 & 2.35 & 6.75 & 3.72 & 1.40\\
$\alpha=90^\circ$$^\dagger$ & 1.55 & -0.48 & 2.62 & 2.30 & 5.29 & -0.39 & 5.21\\
\bottomrule
\end{tabular}
}
  \end{minipage}
\end{table}

\subsection{Distortion sensitivity}\label{sec:sens}
\Cref{tab:ratio} gives the error amplification $e_n(\alpha)/e_n(0^\circ)$ on the
finest mesh, \cref{fig:alpha} the nodal error versus $\alpha$. For a non-polynomial
field neither element is distortion immune: the \PM\ error depends on $\alpha$. For
$\alpha=15^\circ$ and $p\le6$ the \PM\ amplification is $1.5$--$2.7$ (smooth) and
$3.2$--$4.7$ (Lam\'e), against $4.8$--$99$ and $61$--$2900$ for \PP. At $30^\circ$ both
elements are amplified by similar factors ($50$--$2\cdot10^3$ for $p=3$--$6$; \PM\
smaller for some $p$, larger for others). For $\alpha\ge45^\circ$ and $p\ge5$ the \PM\
amplification exceeds the \PP\ amplification by one to ten orders of magnitude. (For
$p\ge7$ the reference errors at $\alpha=0^\circ$ are at round-off level, so these columns
underestimate the sensitivity of \PP\ and overstate that of \PM.)

\begin{table}[t]
  \centering\small
  \caption{Distortion sensitivity on the finest mesh ($N=8$): $e_n(\alpha)/e_n(0^\circ)$,
    smooth field (left) and Lam\'e (right).}
  \label{tab:ratio}
  \begin{minipage}[t]{0.48\linewidth}\centering\resizebox{\linewidth}{!}{\begin{tabular}{lrrrrrrr}
\toprule
& \multicolumn{7}{c}{$p$}\\
\cmidrule(lr){2-8}
 & 2 & 3 & 4 & 5 & 6 & 7 & 8\\
\midrule
\multicolumn{8}{l}{\PM\ (unsymmetric, skew coordinates)}\\
$\alpha=15^\circ$ & 2.25 & 1.54 & 1.93 & 1.8 & 2.56 & 42.8 & 19.3\\
$\alpha=30^\circ$ & 4.24 & 3.15 & 2.2e3 & 50.9 & 27 & 5.5e3 & 7.1e3\\
$\alpha=45^\circ$$^\dagger$ & 204 & 204 & 137 & 5.4e3 & 3.9e4 & 1.9e8 & 1.7e8\\
$\alpha=60^\circ$$^\dagger$ & 144 & 636 & 8.3e3 & 1.3e5 & rej. & rej. & rej.\\
$\alpha=90^\circ$$^\dagger$ & 1.3e3 & 2.4e4 & 9.4e3 & 5.7e3 & 1.2e7 & rej. & rej.\\
\midrule
\multicolumn{8}{l}{\PH\ (unsymmetric, hybrid space)}\\
$\alpha=15^\circ$ & 2.23 & 1.55 & 1.88 & 1.8 & 2.39 & 0.689 & 0.855\\
$\alpha=30^\circ$ & 3.66 & 2.71 & 3.87 & 3.32 & 4.7 & 0.949 & 1.42\\
$\alpha=45^\circ$$^\dagger$ & 5.94 & 6.3 & 6.34 & 8.28 & 8.3 & 30.3 & 1.9e3\\
$\alpha=60^\circ$$^\dagger$ & 8.2e3 & 59.4 & 19.8 & 70 & 452 & rej. & rej.\\
$\alpha=90^\circ$$^\dagger$ & 480 & 187 & 659 & 2.4e3 & 315 & rej. & rej.\\
\midrule
\multicolumn{8}{l}{\PP\ (conventional)}\\
$\alpha=15^\circ$ & 4.84 & 14.4 & 13.2 & 98.7 & 91.1 & 196 & 14.5\\
$\alpha=30^\circ$ & 11.9 & 61.2 & 72.8 & 923 & 1.2e3 & 3.4e3 & 276\\
$\alpha=45^\circ$$^\dagger$ & 22.8 & 221 & 317 & 4.8e3 & 8.6e3 & 7.0e4 & 6.6e3\\
$\alpha=60^\circ$$^\dagger$ & 43.5 & 7.8e3 & 3.3e3 & 4.9e4 & 4.3e5 & 1.3e6 & 4.2e4\\
$\alpha=90^\circ$$^\dagger$ & 654 & 8.1e4 & 1.0e4 & 3.6e5 & 5.0e6 & 3.0e7 & 2.3e6\\
\bottomrule
\end{tabular}
}\end{minipage}\hfill
  \begin{minipage}[t]{0.48\linewidth}\centering\resizebox{\linewidth}{!}{\begin{tabular}{lrrrrrrr}
\toprule
& \multicolumn{7}{c}{$p$}\\
\cmidrule(lr){2-8}
 & 2 & 3 & 4 & 5 & 6 & 7 & 8\\
\midrule
\multicolumn{8}{l}{\PM\ (unsymmetric, skew coordinates)}\\
$\alpha=15^\circ$ & 4.45 & 4.31 & 3.22 & 3.54 & 4.76 & 35.6 & 70.9\\
$\alpha=30^\circ$ & 6.11 & 8.3 & 1.3e3 & 480 & 1.0e3 & 3.9e3 & 1.9e4\\
$\alpha=45^\circ$$^\dagger$ & 253 & 7.9e3 & 1.7e4 & 5.4e4 & 5.3e5 & 4.3e7 & 8.0e9\\
$\alpha=60^\circ$$^\dagger$ & 1.2e3 & 1.8e5 & 6.0e4 & 3.3e6 & rej. & rej. & rej.\\
$\alpha=90^\circ$$^\dagger$ & 3.0e3 & 9.1e3 & 3.9e4 & 1.2e7 & 5.6e7 & rej. & rej.\\
\midrule
\multicolumn{8}{l}{\PH\ (unsymmetric, hybrid space)}\\
$\alpha=15^\circ$ & 7.1 & 6.35 & 12.4 & 33.7 & 88.5 & 42.5 & 1.69\\
$\alpha=30^\circ$ & 23.9 & 15.6 & 45.6 & 111 & 360 & 181 & 7.13\\
$\alpha=45^\circ$$^\dagger$ & 42.9 & 37.5 & 107 & 256 & 1.0e3 & 622 & 1.9e3\\
$\alpha=60^\circ$$^\dagger$ & 811 & 68 & 201 & 523 & 2.6e3 & rej. & rej.\\
$\alpha=90^\circ$$^\dagger$ & 6.6e4 & 991 & 6.3e4 & 6.9e3 & 6.7e4 & rej. & rej.\\
\midrule
\multicolumn{8}{l}{\PP\ (conventional)}\\
$\alpha=15^\circ$ & 61.3 & 132 & 218 & 2.2e3 & 2.3e3 & 311 & 34.6\\
$\alpha=30^\circ$ & 194 & 601 & 1.2e3 & 2.3e4 & 2.9e4 & 5.0e3 & 729\\
$\alpha=45^\circ$$^\dagger$ & 392 & 1.7e3 & 4.3e3 & 7.7e4 & 1.1e5 & 3.3e4 & 5.7e3\\
$\alpha=60^\circ$$^\dagger$ & 619 & 2.4e4 & 2.3e4 & 2.2e5 & 1.1e6 & 1.3e5 & 1.8e5\\
$\alpha=90^\circ$$^\dagger$ & 1.1e3 & 2.3e5 & 6.6e4 & 2.9e6 & 4.2e7 & 3.4e7 & 1.5e6\\
\bottomrule
\end{tabular}
}\end{minipage}
\end{table}

\begin{figure}[t]
  \centering
  \begin{tikzpicture}
\begin{groupplot}[group style={group size=2 by 1, horizontal sep=14mm}, width=6.2cm, height=5.2cm, ymode=log, xtick={0,15,30,45,60,90}, xlabel={$\alpha$ [$^\circ$]}, tick label style={font=\scriptsize}, label style={font=\small}, title style={font=\small}, unbounded coords=jump, every axis plot/.append style={line width=0.7pt, mark size=1.3pt}, grid=major, grid style={gray!20}]
\nextgroupplot[title={Lam\'e, $N=8$}, ylabel={$\en$}]
\addplot[pline=2] table[x=alpha, y=eN2] {\tab/alpha_lame_PM_N8.dat};
\addplot[phyb=2] table[x=alpha, y=eN2] {\tab/alpha_lame_PH_N8.dat};
\addplot[pdash=2] table[x=alpha, y=eN2] {\tab/alpha_lame_PP_N8.dat};
\addplot[pline=4] table[x=alpha, y=eN4] {\tab/alpha_lame_PM_N8.dat};
\addplot[phyb=4] table[x=alpha, y=eN4] {\tab/alpha_lame_PH_N8.dat};
\addplot[pdash=4] table[x=alpha, y=eN4] {\tab/alpha_lame_PP_N8.dat};
\addplot[pline=6] table[x=alpha, y=eN6] {\tab/alpha_lame_PM_N8.dat};
\addplot[phyb=6] table[x=alpha, y=eN6] {\tab/alpha_lame_PH_N8.dat};
\addplot[pdash=6] table[x=alpha, y=eN6] {\tab/alpha_lame_PP_N8.dat};
\addplot[pline=8] table[x=alpha, y=eN8] {\tab/alpha_lame_PM_N8.dat};
\addplot[phyb=8] table[x=alpha, y=eN8] {\tab/alpha_lame_PH_N8.dat};
\addplot[pdash=8] table[x=alpha, y=eN8] {\tab/alpha_lame_PP_N8.dat};
\nextgroupplot[title={smooth, $N=8$}]
\addplot[pline=2] table[x=alpha, y=eN2] {\tab/alpha_smooth_PM_N8.dat};
\addplot[phyb=2] table[x=alpha, y=eN2] {\tab/alpha_smooth_PH_N8.dat};
\addplot[pdash=2] table[x=alpha, y=eN2] {\tab/alpha_smooth_PP_N8.dat};
\addplot[pline=4] table[x=alpha, y=eN4] {\tab/alpha_smooth_PM_N8.dat};
\addplot[phyb=4] table[x=alpha, y=eN4] {\tab/alpha_smooth_PH_N8.dat};
\addplot[pdash=4] table[x=alpha, y=eN4] {\tab/alpha_smooth_PP_N8.dat};
\addplot[pline=6] table[x=alpha, y=eN6] {\tab/alpha_smooth_PM_N8.dat};
\addplot[phyb=6] table[x=alpha, y=eN6] {\tab/alpha_smooth_PH_N8.dat};
\addplot[pdash=6] table[x=alpha, y=eN6] {\tab/alpha_smooth_PP_N8.dat};
\addplot[pline=8] table[x=alpha, y=eN8] {\tab/alpha_smooth_PM_N8.dat};
\addplot[phyb=8] table[x=alpha, y=eN8] {\tab/alpha_smooth_PH_N8.dat};
\addplot[pdash=8] table[x=alpha, y=eN8] {\tab/alpha_smooth_PP_N8.dat};
\end{groupplot}
\matrix[draw=gray!50, fill=white, font=\scriptsize, column sep=1mm, inner sep=1.5pt, anchor=south] at ($(group c1r1.north)!0.5!(group c2r1.north)+(0,9mm)$) { \draw[pline=2] (0,0) -- (4mm,0) node[black, right] {$p=2$}; \pgfmatrixnextcell \draw[pline=4] (0,0) -- (4mm,0) node[black, right] {$p=4$}; \pgfmatrixnextcell \draw[pline=6] (0,0) -- (4mm,0) node[black, right] {$p=6$}; \pgfmatrixnextcell \draw[pline=8] (0,0) -- (4mm,0) node[black, right] {$p=8$}; \pgfmatrixnextcell \draw[black, solid] (0,0) -- (4mm,0) node[black, right] {\PM}; \pgfmatrixnextcell \draw[black, dashdotted] (0,0) -- (4mm,0) node[black, right] {\PH}; \pgfmatrixnextcell \draw[black, dashed] (0,0) -- (4mm,0) node[black, right] {\PP}; \pgfmatrixnextcell \\ };
\end{tikzpicture}

  \caption{Nodal error versus distortion angle on the finest mesh ($N=8$), $p=2,4,6,8$;
    \PM\ solid, \PP\ dashed. Folded elements occur for $\alpha\ge45^\circ$.}
  \label{fig:alpha}
\end{figure}

\begin{figure}[t]
  \centering
  \begin{tikzpicture}
\begin{groupplot}[group style={group size=2 by 1, horizontal sep=14mm}, width=6.2cm, height=5.2cm, ymode=log, xtick={2,...,8}, xlabel={$p$}, tick label style={font=\scriptsize}, label style={font=\small}, title style={font=\small}, unbounded coords=jump, every axis plot/.append style={line width=0.7pt, mark size=1.3pt}, grid=major, grid style={gray!20}]
\nextgroupplot[title={Lam\'e, $N=8$}, ylabel={$e_n$}]
\addplot[color=black, mark=*, mark options={solid, fill=black}] table[x=p, y=eNN8] {\tab/p_lame_PM_a000.dat};
\addplot[color=black, mark=square*, dashdotted, mark options={solid, fill=black}] table[x=p, y=eNN8] {\tab/p_lame_PH_a000.dat};
\addplot[color=black, mark=o, dashed, mark options={solid}] table[x=p, y=eNN8] {\tab/p_lame_PP_a000.dat};
\addplot[color=p3, mark=*, mark options={solid, fill=p3}] table[x=p, y=eNN8] {\tab/p_lame_PM_a015.dat};
\addplot[color=p3, mark=square*, dashdotted, mark options={solid, fill=p3}] table[x=p, y=eNN8] {\tab/p_lame_PH_a015.dat};
\addplot[color=p3, mark=o, dashed, mark options={solid}] table[x=p, y=eNN8] {\tab/p_lame_PP_a015.dat};
\addplot[color=p5, mark=*, mark options={solid, fill=p5}] table[x=p, y=eNN8] {\tab/p_lame_PM_a030.dat};
\addplot[color=p5, mark=square*, dashdotted, mark options={solid, fill=p5}] table[x=p, y=eNN8] {\tab/p_lame_PH_a030.dat};
\addplot[color=p5, mark=o, dashed, mark options={solid}] table[x=p, y=eNN8] {\tab/p_lame_PP_a030.dat};
\addplot[color=p7, mark=*, mark options={solid, fill=p7}] table[x=p, y=eNN8] {\tab/p_lame_PM_a045.dat};
\addplot[color=p7, mark=square*, dashdotted, mark options={solid, fill=p7}] table[x=p, y=eNN8] {\tab/p_lame_PH_a045.dat};
\addplot[color=p7, mark=o, dashed, mark options={solid}] table[x=p, y=eNN8] {\tab/p_lame_PP_a045.dat};
\nextgroupplot[title={smooth MMS, $N=8$}]
\addplot[color=black, mark=*, mark options={solid, fill=black}] table[x=p, y=eNN8] {\tab/p_smooth_PM_a000.dat};
\addplot[color=black, mark=square*, dashdotted, mark options={solid, fill=black}] table[x=p, y=eNN8] {\tab/p_smooth_PH_a000.dat};
\addplot[color=black, mark=o, dashed, mark options={solid}] table[x=p, y=eNN8] {\tab/p_smooth_PP_a000.dat};
\addplot[color=p3, mark=*, mark options={solid, fill=p3}] table[x=p, y=eNN8] {\tab/p_smooth_PM_a015.dat};
\addplot[color=p3, mark=square*, dashdotted, mark options={solid, fill=p3}] table[x=p, y=eNN8] {\tab/p_smooth_PH_a015.dat};
\addplot[color=p3, mark=o, dashed, mark options={solid}] table[x=p, y=eNN8] {\tab/p_smooth_PP_a015.dat};
\addplot[color=p5, mark=*, mark options={solid, fill=p5}] table[x=p, y=eNN8] {\tab/p_smooth_PM_a030.dat};
\addplot[color=p5, mark=square*, dashdotted, mark options={solid, fill=p5}] table[x=p, y=eNN8] {\tab/p_smooth_PH_a030.dat};
\addplot[color=p5, mark=o, dashed, mark options={solid}] table[x=p, y=eNN8] {\tab/p_smooth_PP_a030.dat};
\addplot[color=p7, mark=*, mark options={solid, fill=p7}] table[x=p, y=eNN8] {\tab/p_smooth_PM_a045.dat};
\addplot[color=p7, mark=square*, dashdotted, mark options={solid, fill=p7}] table[x=p, y=eNN8] {\tab/p_smooth_PH_a045.dat};
\addplot[color=p7, mark=o, dashed, mark options={solid}] table[x=p, y=eNN8] {\tab/p_smooth_PP_a045.dat};
\end{groupplot}
\matrix[draw=gray!50, fill=white, font=\scriptsize, column sep=1mm, inner sep=1.5pt, anchor=south] at ($(group c1r1.north)!0.5!(group c2r1.north)+(0,9mm)$) { \draw[color=black, line width=0.7pt] (0,0) -- (4mm,0) node[black, right] {$\alpha=0^\circ$}; \pgfmatrixnextcell \draw[color=p3, line width=0.7pt] (0,0) -- (4mm,0) node[black, right] {$\alpha=15^\circ$}; \pgfmatrixnextcell \draw[color=p5, line width=0.7pt] (0,0) -- (4mm,0) node[black, right] {$\alpha=30^\circ$}; \pgfmatrixnextcell \draw[color=p7, line width=0.7pt] (0,0) -- (4mm,0) node[black, right] {$\alpha=45^\circ$}; \pgfmatrixnextcell \draw[black, solid] (0,0) -- (4mm,0) node[black, right] {\PM}; \pgfmatrixnextcell \draw[black, dashdotted] (0,0) -- (4mm,0) node[black, right] {\PH}; \pgfmatrixnextcell \draw[black, dashed] (0,0) -- (4mm,0) node[black, right] {\PP}; \pgfmatrixnextcell \\ };
\end{tikzpicture}

  \caption{$p$-convergence of the nodal error on the finest mesh ($N=8$), $\alpha=0^\circ,
    15^\circ, 30^\circ, 45^\circ$; \PM\ solid, \PP\ dashed.}
  \label{fig:p}
\end{figure}

\subsection{Mechanism}\label{sec:mechanism}
The diagnostics in \texttt{results/} identify where the loss originates.
\begin{enumerate}
\item \emph{Not the Jacobian, not quadrature.} The loss occurs on meshes with
  $\det\bm J>0$ ($\alpha=30^\circ$) and persists with exact integration
  (\cref{sec:quad,sec:poly}).
\item \emph{Not consistency.} The residual of the exact nodal values,
  $\bm r=\bm K\bm u_I-\bm F$, is \emph{smaller} for \PM\ than for \PP\ in all eleven cases
  tested (\texttt{residual\_smooth.txt}; e.g.\ $\alpha=30^\circ$, $N=8$, $p=4$:
  $3.0\cdot10^{-5}$ versus $3.7\cdot10^{-4}$), whereas the \PM\ nodal error is $9$ times
  larger ($2.8\cdot10^{-5}$ versus $3.0\cdot10^{-6}$). The ratio of error to residual of
  \PM\ grows under refinement ($0.05$ at $N=4$, $0.9$ at $N=8$).
\item \emph{Stability of the unsymmetric operator.} With $\bm u_h-\bm u_I=-\bm K^{-1}\bm r=
  -\sum_k\bm v_k(\bm u_k^{\mathsf T}\bm r)/\sigma_k$ (singular triplets of $\bm K$;
  \texttt{spectrum\_smooth.txt}, smooth field), the error of \PP\ and of \PM\ for
  $\alpha\le15^\circ$ is spread over the spectrum (largest single contribution $2$--$8\,\%$ of
  $\|\bm u_h-\bm u_I\|^2$, from mid-spectrum modes). For $\alpha\ge30^\circ$, $p=4$, \PM\ has
  singular values below the smallest singular value of \PP\ (down to $0.17\,\sigma_1^{\PP}$),
  and in four of the five cases tested a single one of the lowest six triplets carries
  $21$--$68\,\%$ of the error (the fifth, $\alpha=30^\circ$, $N=4$, $p=6$: triplet $k=28$,
  $9\,\%$). The
  residual of the non-conforming Petrov--Galerkin scheme is not orthogonal to these
  low modes, unlike the Galerkin residual of \PP, so they are excited and amplified.
\item \emph{Ill-conditioned metric interpolation.} The Lebesgue constant
  $\Lambda_M=\max_{\bx\in\Omega^e}\sum_i|M_i^e(\bx)|$ of the metric functions (basis
  independent) equals the GLL value on affine elements, but grows by orders of magnitude
  with $\alpha$ and $p$ (\cref{tab:leb}): $\Lambda_M\le163$ for $\alpha=15^\circ$, $p\le6$,
  but $1.1\cdot10^3$ ($\alpha=30^\circ$, $p=4$), $9\cdot10^5$ ($45^\circ$, $p=5$) and
  $5\cdot10^{12}$ ($60^\circ$, $p=8$). The element condition numbers of \PM\ exceed those of
  \PP\ by similar factors (\cref{tab:cond}). The images of the GLL nodes of strongly
  non-affine elements come close to configurations that are not unisolvent for $Q_p$ in
  skew coordinates (``interpolation failure'' in \citep{Rajendran2010technique}); the
  metric trial functions then have large gradients, the bilinear form loses its stability
  constant, and round-off is amplified (\cref{sec:poly}). All cases of the benign regime
  of \cref{sec:conv} have $\Lambda_M\le2.5\cdot10^2$, all degraded cases
  $\Lambda_M\ge10^3$; the converse does not hold strictly ($\alpha=30^\circ$, $p=6$:
  $\Lambda_M=1.1\cdot10^4$, yet \PM\ is $500$ times more accurate than \PP\ at $N=8$), so
  $\Lambda_M$ is a necessary rather than a sufficient indicator. The maximum over the mesh does not decrease under
  refinement because the distortion is self-similar; it varies with $N$ because the
  element shapes near the curved boundaries change.
\end{enumerate}

\begin{table}[t]
  \centering\small
  \begin{minipage}[t]{0.48\linewidth}\centering
    \captionof{table}{Lebesgue constant $\Lambda_M$ of the metric interpolation, maximum over
      the elements, $N=8$ (GLL: affine element).}\label{tab:leb}
    \resizebox{\linewidth}{!}{\begin{tabular}{lrrrrrrr}
\toprule
& \multicolumn{7}{c}{$p$}\\
\cmidrule(lr){2-8}
 & 2 & 3 & 4 & 5 & 6 & 7 & 8\\
\midrule
\multicolumn{8}{l}{\PM\ (unsymmetric, skew coordinates)}\\
$\alpha=0^\circ$ & $1.56$ & $2.22$ & $2.68$ & $3.02$ & $3.42$ & $4.60$ & $5.35$\\
$\alpha=15^\circ$ & $2.53$ & $5.98$ & $15.5$ & $37.2$ & $130$ & $1.2\cdot10^{4}$ & $1.3\cdot10^{4}$\\
$\alpha=30^\circ$ & $9.95$ & $240$ & $1.1\cdot10^{3}$ & $1.6\cdot10^{3}$ & $6.2\cdot10^{3}$ & $4.2\cdot10^{6}$ & $1.2\cdot10^{6}$\\
$\alpha=45^\circ$$^\dagger$ & $186$ & $1.5\cdot10^{3}$ & $1.3\cdot10^{4}$ & $5.0\cdot10^{5}$ & $8.3\cdot10^{5}$ & $1.4\cdot10^{6}$ & $1.3\cdot10^{8}$\\
$\alpha=60^\circ$$^\dagger$ & $290$ & $1.1\cdot10^{4}$ & $8.5\cdot10^{4}$ & $2.8\cdot10^{5}$ & rej. & rej. & rej.\\
$\alpha=90^\circ$$^\dagger$ & $61.2$ & $3.4\cdot10^{3}$ & $8.3\cdot10^{5}$ & $2.2\cdot10^{5}$ & $3.3\cdot10^{5}$ & rej. & rej.\\
\midrule
\multicolumn{8}{l}{\PH\ (unsymmetric, hybrid space)}\\
$\alpha=0^\circ$ & $1.58$ & $2.22$ & $2.69$ & $3.00$ & $3.41$ & $3.47$ & $4.15$\\
$\alpha=15^\circ$ & $1.76$ & $2.40$ & $2.80$ & $3.10$ & $3.47$ & $4.47$ & $4.18$\\
$\alpha=30^\circ$ & $2.00$ & $4.77$ & $7.52$ & $16.0$ & $24.9$ & $74.5$ & $94.6$\\
$\alpha=45^\circ$$^\dagger$ & $2.78$ & $10.0$ & $21.2$ & $46.7$ & $91.0$ & $258$ & $546$\\
$\alpha=60^\circ$$^\dagger$ & $4.26$ & $18.9$ & $53.3$ & $193$ & $833$ & rej. & rej.\\
$\alpha=90^\circ$$^\dagger$ & $389$ & $2.9\cdot10^{3}$ & $4.8\cdot10^{3}$ & $1.7\cdot10^{3}$ & $2.2\cdot10^{3}$ & rej. & rej.\\
\midrule
GLL & $1.56$ & $2.22$ & $2.68$ & $3.00$ & $3.41$ & $3.47$ & $4.15$\\
\bottomrule
\end{tabular}
}
  \end{minipage}\hfill
  \begin{minipage}[t]{0.48\linewidth}\centering
    \captionof{table}{Ratio of the maximum element condition numbers
      $\sigma_{\max}/\sigma_{2n-3}$, \PM/\PP, $N=8$.}\label{tab:cond}
    \resizebox{\linewidth}{!}{\begin{tabular}{lrrrrrrr}
\toprule
& \multicolumn{7}{c}{$p$}\\
\cmidrule(lr){2-8}
 & 2 & 3 & 4 & 5 & 6 & 7 & 8\\
\midrule
$\alpha=0^\circ$ & 0.998 & 1.01 & 0.996 & 1 & 0.997 & 1.01 & 0.998\\
$\alpha=15^\circ$ & 1.54 & 0.937 & 1.19 & 1.61 & 3.25 & 182 & 128\\
$\alpha=30^\circ$ & 60 & 11.5 & 145 & 623 & 90.3 & 1.4e5 & 2.1e6\\
$\alpha=45^\circ$$^\dagger$ & 44.4 & 12.8 & 148 & 7.2e4 & 1.1e5 & 6.0e6 & 2.4e6\\
$\alpha=60^\circ$$^\dagger$ & 10.4 & 10.7 & 104 & 7.2e3 & 2.2e5 & 3.6e7 & 3.5e10\\
$\alpha=90^\circ$$^\dagger$ & 9.95 & 198 & 23.2 & 5.7e3 & 3.7e4 & 1.6e5 & 2.8e7\\
\bottomrule
\end{tabular}
}
  \end{minipage}
\end{table}


\subsection{Quadrature}\label{sec:quad}
The \PM\ integrand \eqref{eq:Kpm} is a polynomial of degree $(2p-1)(q+1)$ per direction
for a geometry of degree $q$; exact integration needs $n_q\ge p^2+p/2$ points on the
elements with exact arcs ($q=p$) and $n_q\ge3p-1$ on the interior elements ($q=2$).
The default rule ($n_q=2p$) is therefore not exact. Repeating selected cases with
$n_q=2p,3p,4p,6p$ (\texttt{results/raw/quadrature.csv}; $\alpha=0^\circ,30^\circ$, $N=1,2$,
$p=5,8$) changes $\eE$ of the non-polynomial fields by at most $3\cdot10^{-7}$ (relative) for
\PP\ and $0.5\,\%$ for \PM; the single exception (\PM, smooth field, $\alpha=30^\circ$,
$N=2$, $p=8$: $39\,\%$) has an error only six times above the round-off level of that
mesh. For the polynomial field (\PM, $p\ge5$), whose error is pure round-off, the value
changes by up to a factor $8$. The conclusions are therefore not affected by quadrature. On folded meshes
the results change erratically with $n_q$ (e.g.\ polynomial field, $\alpha=60^\circ$,
$N=2$, $p=8$: $\eE=0.78$, $0.35$, $0.10$, $0.38$), consistent with the round-off
amplification of \cref{sec:poly}, which persists with exact integration.

\begin{figure}[t]
  \centering
  \begin{tikzpicture}
\begin{groupplot}[group style={group size=6 by 3, horizontal sep=1.5mm, vertical sep=8mm, yticklabels at=edge left, xticklabels at=edge bottom}, conv, width=3.90cm, height=4.4cm, ymin=1e-15, ymax=1e2]
\nextgroupplot[title={\PM, $\alpha=0^\circ$}, ylabel={$e_n$}]
\addplot[pline=2] table[x=h, y=eN2] {\tab/h_poly_PM_a000.dat};
\addplot[pline=3] table[x=h, y=eN3] {\tab/h_poly_PM_a000.dat};
\addplot[pline=4] table[x=h, y=eN4] {\tab/h_poly_PM_a000.dat};
\addplot[pline=5] table[x=h, y=eN5] {\tab/h_poly_PM_a000.dat};
\addplot[pline=6] table[x=h, y=eN6] {\tab/h_poly_PM_a000.dat};
\addplot[pline=7] table[x=h, y=eN7] {\tab/h_poly_PM_a000.dat};
\addplot[pline=8] table[x=h, y=eN8] {\tab/h_poly_PM_a000.dat};
\nextgroupplot[title={\PM, $\alpha=15^\circ$}]
\addplot[pline=2] table[x=h, y=eN2] {\tab/h_poly_PM_a015.dat};
\addplot[pline=3] table[x=h, y=eN3] {\tab/h_poly_PM_a015.dat};
\addplot[pline=4] table[x=h, y=eN4] {\tab/h_poly_PM_a015.dat};
\addplot[pline=5] table[x=h, y=eN5] {\tab/h_poly_PM_a015.dat};
\addplot[pline=6] table[x=h, y=eN6] {\tab/h_poly_PM_a015.dat};
\addplot[pline=7] table[x=h, y=eN7] {\tab/h_poly_PM_a015.dat};
\addplot[pline=8] table[x=h, y=eN8] {\tab/h_poly_PM_a015.dat};
\nextgroupplot[title={\PM, $\alpha=30^\circ$}]
\addplot[pline=2] table[x=h, y=eN2] {\tab/h_poly_PM_a030.dat};
\addplot[pline=3] table[x=h, y=eN3] {\tab/h_poly_PM_a030.dat};
\addplot[pline=4] table[x=h, y=eN4] {\tab/h_poly_PM_a030.dat};
\addplot[pline=5] table[x=h, y=eN5] {\tab/h_poly_PM_a030.dat};
\addplot[pline=6] table[x=h, y=eN6] {\tab/h_poly_PM_a030.dat};
\addplot[pline=7] table[x=h, y=eN7] {\tab/h_poly_PM_a030.dat};
\addplot[pline=8] table[x=h, y=eN8] {\tab/h_poly_PM_a030.dat};
\nextgroupplot[title={\PM, $\alpha=45^\circ$}]
\addplot[pline=2] table[x=h, y=eN2] {\tab/h_poly_PM_a045.dat};
\addplot[pline=3] table[x=h, y=eN3] {\tab/h_poly_PM_a045.dat};
\addplot[pline=4] table[x=h, y=eN4] {\tab/h_poly_PM_a045.dat};
\addplot[pline=5] table[x=h, y=eN5] {\tab/h_poly_PM_a045.dat};
\addplot[pline=6] table[x=h, y=eN6] {\tab/h_poly_PM_a045.dat};
\addplot[pline=7] table[x=h, y=eN7] {\tab/h_poly_PM_a045.dat};
\addplot[pline=8] table[x=h, y=eN8] {\tab/h_poly_PM_a045.dat};
\nextgroupplot[title={\PM, $\alpha=60^\circ$}]
\addplot[pline=2] table[x=h, y=eN2] {\tab/h_poly_PM_a060.dat};
\addplot[pline=3] table[x=h, y=eN3] {\tab/h_poly_PM_a060.dat};
\addplot[pline=4] table[x=h, y=eN4] {\tab/h_poly_PM_a060.dat};
\addplot[pline=5] table[x=h, y=eN5] {\tab/h_poly_PM_a060.dat};
\addplot[pline=6] table[x=h, y=eN6] {\tab/h_poly_PM_a060.dat};
\addplot[pline=7] table[x=h, y=eN7] {\tab/h_poly_PM_a060.dat};
\addplot[pline=8] table[x=h, y=eN8] {\tab/h_poly_PM_a060.dat};
\nextgroupplot[title={\PM, $\alpha=90^\circ$}]
\addplot[pline=2] table[x=h, y=eN2] {\tab/h_poly_PM_a090.dat};
\addplot[pline=3] table[x=h, y=eN3] {\tab/h_poly_PM_a090.dat};
\addplot[pline=4] table[x=h, y=eN4] {\tab/h_poly_PM_a090.dat};
\addplot[pline=5] table[x=h, y=eN5] {\tab/h_poly_PM_a090.dat};
\addplot[pline=6] table[x=h, y=eN6] {\tab/h_poly_PM_a090.dat};
\addplot[pline=7] table[x=h, y=eN7] {\tab/h_poly_PM_a090.dat};
\addplot[pline=8] table[x=h, y=eN8] {\tab/h_poly_PM_a090.dat};
\nextgroupplot[title={\PH, $\alpha=0^\circ$}, ylabel={$e_n$}]
\addplot[pline=2] table[x=h, y=eN2] {\tab/h_poly_PH_a000.dat};
\addplot[pline=3] table[x=h, y=eN3] {\tab/h_poly_PH_a000.dat};
\addplot[pline=4] table[x=h, y=eN4] {\tab/h_poly_PH_a000.dat};
\addplot[pline=5] table[x=h, y=eN5] {\tab/h_poly_PH_a000.dat};
\addplot[pline=6] table[x=h, y=eN6] {\tab/h_poly_PH_a000.dat};
\addplot[pline=7] table[x=h, y=eN7] {\tab/h_poly_PH_a000.dat};
\addplot[pline=8] table[x=h, y=eN8] {\tab/h_poly_PH_a000.dat};
\nextgroupplot[title={\PH, $\alpha=15^\circ$}]
\addplot[pline=2] table[x=h, y=eN2] {\tab/h_poly_PH_a015.dat};
\addplot[pline=3] table[x=h, y=eN3] {\tab/h_poly_PH_a015.dat};
\addplot[pline=4] table[x=h, y=eN4] {\tab/h_poly_PH_a015.dat};
\addplot[pline=5] table[x=h, y=eN5] {\tab/h_poly_PH_a015.dat};
\addplot[pline=6] table[x=h, y=eN6] {\tab/h_poly_PH_a015.dat};
\addplot[pline=7] table[x=h, y=eN7] {\tab/h_poly_PH_a015.dat};
\addplot[pline=8] table[x=h, y=eN8] {\tab/h_poly_PH_a015.dat};
\nextgroupplot[title={\PH, $\alpha=30^\circ$}]
\addplot[pline=2] table[x=h, y=eN2] {\tab/h_poly_PH_a030.dat};
\addplot[pline=3] table[x=h, y=eN3] {\tab/h_poly_PH_a030.dat};
\addplot[pline=4] table[x=h, y=eN4] {\tab/h_poly_PH_a030.dat};
\addplot[pline=5] table[x=h, y=eN5] {\tab/h_poly_PH_a030.dat};
\addplot[pline=6] table[x=h, y=eN6] {\tab/h_poly_PH_a030.dat};
\addplot[pline=7] table[x=h, y=eN7] {\tab/h_poly_PH_a030.dat};
\addplot[pline=8] table[x=h, y=eN8] {\tab/h_poly_PH_a030.dat};
\nextgroupplot[title={\PH, $\alpha=45^\circ$}]
\addplot[pline=2] table[x=h, y=eN2] {\tab/h_poly_PH_a045.dat};
\addplot[pline=3] table[x=h, y=eN3] {\tab/h_poly_PH_a045.dat};
\addplot[pline=4] table[x=h, y=eN4] {\tab/h_poly_PH_a045.dat};
\addplot[pline=5] table[x=h, y=eN5] {\tab/h_poly_PH_a045.dat};
\addplot[pline=6] table[x=h, y=eN6] {\tab/h_poly_PH_a045.dat};
\addplot[pline=7] table[x=h, y=eN7] {\tab/h_poly_PH_a045.dat};
\addplot[pline=8] table[x=h, y=eN8] {\tab/h_poly_PH_a045.dat};
\nextgroupplot[title={\PH, $\alpha=60^\circ$}]
\addplot[pline=2] table[x=h, y=eN2] {\tab/h_poly_PH_a060.dat};
\addplot[pline=3] table[x=h, y=eN3] {\tab/h_poly_PH_a060.dat};
\addplot[pline=4] table[x=h, y=eN4] {\tab/h_poly_PH_a060.dat};
\addplot[pline=5] table[x=h, y=eN5] {\tab/h_poly_PH_a060.dat};
\addplot[pline=6] table[x=h, y=eN6] {\tab/h_poly_PH_a060.dat};
\addplot[pline=7] table[x=h, y=eN7] {\tab/h_poly_PH_a060.dat};
\addplot[pline=8] table[x=h, y=eN8] {\tab/h_poly_PH_a060.dat};
\nextgroupplot[title={\PH, $\alpha=90^\circ$}]
\addplot[pline=2] table[x=h, y=eN2] {\tab/h_poly_PH_a090.dat};
\addplot[pline=3] table[x=h, y=eN3] {\tab/h_poly_PH_a090.dat};
\addplot[pline=4] table[x=h, y=eN4] {\tab/h_poly_PH_a090.dat};
\addplot[pline=5] table[x=h, y=eN5] {\tab/h_poly_PH_a090.dat};
\addplot[pline=6] table[x=h, y=eN6] {\tab/h_poly_PH_a090.dat};
\addplot[pline=7] table[x=h, y=eN7] {\tab/h_poly_PH_a090.dat};
\addplot[pline=8] table[x=h, y=eN8] {\tab/h_poly_PH_a090.dat};
\nextgroupplot[title={\PP, $\alpha=0^\circ$}, ylabel={$e_n$}]
\addplot[pline=2] table[x=h, y=eN2] {\tab/h_poly_PP_a000.dat};
\addplot[pline=3] table[x=h, y=eN3] {\tab/h_poly_PP_a000.dat};
\addplot[pline=4] table[x=h, y=eN4] {\tab/h_poly_PP_a000.dat};
\addplot[pline=5] table[x=h, y=eN5] {\tab/h_poly_PP_a000.dat};
\addplot[pline=6] table[x=h, y=eN6] {\tab/h_poly_PP_a000.dat};
\addplot[pline=7] table[x=h, y=eN7] {\tab/h_poly_PP_a000.dat};
\addplot[pline=8] table[x=h, y=eN8] {\tab/h_poly_PP_a000.dat};
\nextgroupplot[title={\PP, $\alpha=15^\circ$}]
\addplot[pline=2] table[x=h, y=eN2] {\tab/h_poly_PP_a015.dat};
\addplot[pline=3] table[x=h, y=eN3] {\tab/h_poly_PP_a015.dat};
\addplot[pline=4] table[x=h, y=eN4] {\tab/h_poly_PP_a015.dat};
\addplot[pline=5] table[x=h, y=eN5] {\tab/h_poly_PP_a015.dat};
\addplot[pline=6] table[x=h, y=eN6] {\tab/h_poly_PP_a015.dat};
\addplot[pline=7] table[x=h, y=eN7] {\tab/h_poly_PP_a015.dat};
\addplot[pline=8] table[x=h, y=eN8] {\tab/h_poly_PP_a015.dat};
\nextgroupplot[title={\PP, $\alpha=30^\circ$}]
\addplot[pline=2] table[x=h, y=eN2] {\tab/h_poly_PP_a030.dat};
\addplot[pline=3] table[x=h, y=eN3] {\tab/h_poly_PP_a030.dat};
\addplot[pline=4] table[x=h, y=eN4] {\tab/h_poly_PP_a030.dat};
\addplot[pline=5] table[x=h, y=eN5] {\tab/h_poly_PP_a030.dat};
\addplot[pline=6] table[x=h, y=eN6] {\tab/h_poly_PP_a030.dat};
\addplot[pline=7] table[x=h, y=eN7] {\tab/h_poly_PP_a030.dat};
\addplot[pline=8] table[x=h, y=eN8] {\tab/h_poly_PP_a030.dat};
\nextgroupplot[title={\PP, $\alpha=45^\circ$}]
\addplot[pline=2] table[x=h, y=eN2] {\tab/h_poly_PP_a045.dat};
\addplot[pline=3] table[x=h, y=eN3] {\tab/h_poly_PP_a045.dat};
\addplot[pline=4] table[x=h, y=eN4] {\tab/h_poly_PP_a045.dat};
\addplot[pline=5] table[x=h, y=eN5] {\tab/h_poly_PP_a045.dat};
\addplot[pline=6] table[x=h, y=eN6] {\tab/h_poly_PP_a045.dat};
\addplot[pline=7] table[x=h, y=eN7] {\tab/h_poly_PP_a045.dat};
\addplot[pline=8] table[x=h, y=eN8] {\tab/h_poly_PP_a045.dat};
\nextgroupplot[title={\PP, $\alpha=60^\circ$}]
\addplot[pline=2] table[x=h, y=eN2] {\tab/h_poly_PP_a060.dat};
\addplot[pline=3] table[x=h, y=eN3] {\tab/h_poly_PP_a060.dat};
\addplot[pline=4] table[x=h, y=eN4] {\tab/h_poly_PP_a060.dat};
\addplot[pline=5] table[x=h, y=eN5] {\tab/h_poly_PP_a060.dat};
\addplot[pline=6] table[x=h, y=eN6] {\tab/h_poly_PP_a060.dat};
\addplot[pline=7] table[x=h, y=eN7] {\tab/h_poly_PP_a060.dat};
\addplot[pline=8] table[x=h, y=eN8] {\tab/h_poly_PP_a060.dat};
\nextgroupplot[title={\PP, $\alpha=90^\circ$}]
\addplot[pline=2] table[x=h, y=eN2] {\tab/h_poly_PP_a090.dat};
\addplot[pline=3] table[x=h, y=eN3] {\tab/h_poly_PP_a090.dat};
\addplot[pline=4] table[x=h, y=eN4] {\tab/h_poly_PP_a090.dat};
\addplot[pline=5] table[x=h, y=eN5] {\tab/h_poly_PP_a090.dat};
\addplot[pline=6] table[x=h, y=eN6] {\tab/h_poly_PP_a090.dat};
\addplot[pline=7] table[x=h, y=eN7] {\tab/h_poly_PP_a090.dat};
\addplot[pline=8] table[x=h, y=eN8] {\tab/h_poly_PP_a090.dat};
\end{groupplot}
\node[font=\small] at ($(group c1r3.south)!0.5!(group c6r3.south)+(0,-8mm)$) {element size $h\propto 1/N$};
\matrix[draw=gray!50, fill=white, font=\scriptsize, column sep=1mm, inner sep=1.5pt, anchor=south] at ($(group c1r1.north)!0.5!(group c6r1.north)+(0,7mm)$) { \draw[pline=2] (0,0) -- (4mm,0) node[black, right] {$p=2$}; \pgfmatrixnextcell \draw[pline=3] (0,0) -- (4mm,0) node[black, right] {$p=3$}; \pgfmatrixnextcell \draw[pline=4] (0,0) -- (4mm,0) node[black, right] {$p=4$}; \pgfmatrixnextcell \draw[pline=5] (0,0) -- (4mm,0) node[black, right] {$p=5$}; \pgfmatrixnextcell \draw[pline=6] (0,0) -- (4mm,0) node[black, right] {$p=6$}; \pgfmatrixnextcell \draw[pline=7] (0,0) -- (4mm,0) node[black, right] {$p=7$}; \pgfmatrixnextcell \draw[pline=8] (0,0) -- (4mm,0) node[black, right] {$p=8$}; \pgfmatrixnextcell  \\ };

\end{tikzpicture}

  \caption{Polynomial field (degree 5): nodal error versus element size, \PM\ (top) and
    \PP\ (bottom). \PM\ reproduces the field for $p\ge5$ up to round-off on all meshes
    without folded elements.}
  \label{fig:h-poly}
\end{figure}
